\cotefile{157-4}
\scanpage{180}
\mspage{54}
% archive sheet number: 179

\unsure{prouver} que $M$ est stable par $\varinjlim$ grandes devant $\unsure{\pi|\Pi}$ sans plus, et je doute maintenant que ça résulte des seules hypothèses faites ici, qui ne concernent que les systèmes inductifs où les morphismes de transition sont mono.

Voici, je crois, un contre-exemple. Prenons

\[
\begin{aligned}
 M={}&\text{ensemble des objets de $\mathcal M$ isomorphes}\\
 &\text{à un objet \unsure{fixe|fixé} $\deleted{\unsure{X}}X_0\in\mathrm{Filt}^{\Pi}_{\deleted{\unsure{\pi}}<}$,}\\
 &\text{tel que $X_0$ n'ait pas d'\unsure{endomorphismes|endomorphisme}}\\
 &\text{qui soit mono que l'identité.}
\end{aligned}\tag*{(5.4.)}
\]

Alors la condition 5.3. est satisfaite. Car \deleted{\unsure{a)}}
a) si $X\in M$, i.e. $X\simeq X_0$, donc $X\in\mathrm{Filt}^{\Pi}_{\deleted{\unsure{\pi|<}}}$, \deleted{\unsure{donc}}
$\mathrm{Ss}_{\pi}(X)$ admet \deleted{\unsure{un}} \inserted{$X$ comme} objet final et $J$, \deleted{étant} cofinal, contient \deleted{$X$,} \inserted{cet \unsure{élément|objet}}, donc $J_M$ (qui le contient aussi) est cofinal. Et b) si $J_M$ est cofinal, \deleted{\illegible{1}} l'hypothèse (5.4.) implique que les morphismes \deleted{de} transition \deleted{\unsure{dans}} entre les $Z_\alpha$ ($\alpha\in J_M$) est formé d'isomorphismes, \deleted{donc} la limite inductive $X$ des $Z_\alpha$ en question est isom. à chacun d'eux, \deleted{donc} $X\deleted{\in M.}\simeq X_0$ i.e. $X\in M$.

\unsure{Pourtant|Par contre}, $M$ n'est $\pi$-accessible que s'il est stable par $\varinjlim$ grandes devant $\pi$, ce qui implique qu'il est stable par \uline{facteurs directs}. \deleted{\unsure{Donc Mais l'hypothèse}}

\scanpage{181}
\mspage{55}
% archive sheet number: 180

(5.4) n'exclut nullement que $X_0$ admette un projecteur $p$ ($p^2=p$) non trivial. Celui-ci aura une image $X'$ \deleted{dans} $\mathcal M$, qui sera facteur direct de $X_0$ non isom. à $X_0$ i.e. non dans $\deleted{\unsure{M}}M$.

Ces difficultés disparaissent si on suppose que $M^{\unsure{\pi|\Pi}}$ est stable par $\unsure{\underline{\lim}|\varinjlim}$ filtrantes indexées par des ens. \deleted{\unsure{ord}} ordonnés $I$ de cardinal $\leq\pi$. En effet, dans ce cas, \deleted{\unsure{le sous}} dans la démonstration précédente, le sous-ensemble $I_M$ est stable par bornes supérieures (dans $I$) des parties $J$ de $I_M$ filtrantes et de cardinal $\leq\pi$. Donc le syst. inductif $(\unsure{X|X'}_i)_{i\in I_M}$ est encore $\pi$-adapté, et on peut lui appliquer le Lemme 5.7.) pour trouver un sous-ensemble cofinal $I'_M\subset I_M$, tel que les $X_i\unsure{\longrightarrow|-} X$ pour $i\in I'_M$ soient monos. L'image $J$ de $I'_M$ dans $\mathrm{Ss}_{\pi}(X)$ est alors une partie de $\mathrm{Ss}_{\pi}(X)$ telle que $\varinjlim_{\alpha\in J} Z_\alpha=X$, donc (par l'hyp. faite sur $\pi$) cofinale dans $\mathrm{Ss}_{\pi}(X)$ (5.4.), avec $Z_\alpha\in M^\pi$ \deleted{\unsure{qq}} pour tt $\alpha\in J$. Donc par (5.3.) on a $X\in M$, qed.

\scanpage{182}
\mspage{56}
% archive sheet number: 181

Il reste à prouver le Lemme 5.7. \deleted{\unsure{Avec}} Avec les hyp. \inserted{et} notations du lemme, soit donné $i_0\in I$, il faut trouver $i'\in I'$ tel que $i'\geq i_0$. Considérons l'image $X'_i$ de $X_i$ dans $X$, et l'inclusion

\[
\underset{\substack{\rotatebox{-90}{$\in$}\\\mathrm{Filt}^{\unsure{\pi|\Pi}}}}{X'_{i_0}}\longrightarrow X\simeq\varinjlim X_i
\]

Comme $X'_{i_0}\in\mathrm{Filt}^\pi$, donc \inserted{\unsure{est}} $\pi$-accessible, il existe un $i_1>i_0$ tel que $X'_{i_0}\unsure{\longrightarrow|-} X$ se factorise par

\[
X'_{i_0}\longrightarrow X_{i_1}.
\]

Appliquant le $=$ \unsure{résultat} à $X_{i_1},X'_{i_1}$ etc., on trouve une suite \inserted{infinie} de morphismes

\diagram{\begin{tikzcd}[ampersand replacement=\&,column sep=small,row sep=small]
 X_{i_0}\arrow[dr] \& \& X_{i_1}\arrow[dr] \& \& X_{i_2}\arrow[dr,no head] \& \& \\
 \& X'_{i_0}\arrow[ur] \& \& X'_{i_1}\arrow[ur] \& \& X'_{\unsure{i_2|j_2}} \& \cdots
\end{tikzcd}}

telle que les morphismes composés de deux en deux soient, \deleted{\unsure{les}} sur la ligne supérieure, les transitions

\[
X_{i_{\unsure{m|n}}}\unsure{-|\longrightarrow} X_{i_p}\qquad(i_{\unsure{m|p}}>i_{\unsure{n|m}})
\]

et sur la ligne inférieure, les inclusions

\[
 X'_{i_m}\lhook\joinrel\longrightarrow X'_{i_n}\quad(\unsure{\lhook\joinrel\longrightarrow|\lhook\joinrel-} X).
\]

Soit $\deleted{\unsure{X}}\,i_\omega=\underset{n}{\mathrm{Sup}}\,i_n$ dans $I$, \deleted{\unsure{donc}}

\[
 X_{i_\omega}=\varinjlim_n X^{\deleted{\unsure{\prime}}}_{i_n}\simeq\varinjlim_n X'_{i_n}\simeq X'_{i_\omega},
\]

donc $X_{i_\omega}\longrightarrow X$ est mono.

\scanpage{183}
\mspage{57}
% archive sheet number: 182

La \unsure{dernière|deuxième} assertion de 5.7. résulte aussitôt du fait qu'une limite inductive filtrante de monos dans $\mathcal M$ est mono. (condition e) p. 46)

On peut préciser la prop. 5.6. obtenue ainsi :

\uline{Proposition 5.7.} \deleted{\unsure{Soient}} Soit $\mathcal M$ une catégorie satisfaisant les conditions \unsure{a) b)} c) d) (p. 44,45\unsure{,|-} et e) (p. 46). Soit $M$ une partie de $\mathrm{Ob}\,\mathcal M$, stable par isomorphie, \deleted{les conditions \unsure{suivantes} \illegible{1}} \inserted{$\unsure{\pi|\Pi}$ un cardinal tel} que les objets de $\mathrm{Filt}^{\Pi}$ soient $\pi$-accessibles. Les conditions suivantes sont équivalentes :

a) $M$ est $\pi$-accessible, et stable par $\varinjlim$ filtrantes

b) $M$ est $\pi$ accessible, et $\inserted{M\cap\mathrm{Filt}^{\Pi}}$ \deleted{\unsure{est}} est stable par $\underline{\lim}$ filtrantes \deleted{\unsure{de}} cardinal $\leq\pi$.

% slanted
\margin{\deleted{Condition de \unsure{$\pi$-compacité} qui peut se vérifier \unsure{maintenant} en \unsure{5.1)} p. 50, cf en bas..} \unsure{$\to$}}

c) Il existe une partie $M^\pi$ de $\mathrm{Filt}^{\Pi}$, stable par $\underline{\lim}$ filtrantes de cardinal $\leq\pi$, \deleted{\unsure{et pour}} et telle que pour tout $X\in\mathrm{Ob}\,\mathcal M$, \deleted{\unsure{toute} partie $J$ cofinale de $\mathrm{Ss}_\pi(X)$ (stable par bornes supérieures de parties filtrantes de cardinal $\leq\pi$)} on ait \deleted{\unsure{équivalence}} $X\in M$ sss l'ens $J_M$ des $Z_\alpha\in J$ tels que $Z_\alpha\in M^\pi$ \deleted{\unsure{soit}} est cofinal dans $J$.

(De plus, sous la condition c) on a

\scanpage{184}
\mspage{58}
% archive sheet number: 183

\[
M^{\unsure{\pi|\Pi}}=M\cap\mathrm{Filt}^{\Pi}.)
\]

\inserted{c$'$) Il existe une partie $M_0$ de $\mathrm{Filt}^{\Pi}$, telle que $M$ soit la plus petite partie de $\mathcal M$ stable par limites ind. filtrantes. (De plus, on peut alors prendre $M_0$ \uline{petit})}

% slanted
\margin{\uline{NB} On
\diagram{\begin{tikzcd}[ampersand replacement=\&,column sep=tiny,row sep=small]
 a\arrow[r,Leftrightarrow]\arrow[d,Rightarrow] \& b\arrow[d,Rightarrow] \& \deleted{\unsure{\Rightarrow c}}\\
 c')\arrow[r,Rightarrow,"{\textcircled{?}}"] \& c) \& \deleted{\unsure{\Rightarrow}}
\end{tikzcd}}
\unsure{sans} condition sur $\pi$ (sauf $\pi\geq\pi_0$)}

\uline{Corollaire 5.8.} Les applications $M\longmapsto M^\pi=M\cap\mathrm{Filt}^{\Pi}$,

$M^\pi\longmapsto$ l'ensemble des $X\in\mathrm{Ob}\,\mathcal M$ tels que l'ens $\mathrm{Ss}^{M^\pi}_{\pi}(X)$ des $Z_\alpha\in\mathrm{Ss}_{\pi}(X)$ qui sont dans $M^\pi$ soit cofinal
\inserted{($=$ \deleted{\unsure{le}} \deleted{plus} petit sous-ens. de $\mathrm{Ob}\,\mathcal M$ stable par $\varinjlim$ filtrantes)}

\unsure{établissent|et définissent} des bijections, inverses l'une de l'autre, entre l'ens des $M\subset\mathrm{Ob}\,\mathcal M$ satisfaisant les conditions \inserted{équivalentes} de la prop. 5.7, et l'ensemble des parties $M^\pi$ de $\mathrm{Filt}^{\Pi}$ stables par $\varinjlim$ filtrantes de cardinal $\leq\pi$.

\uline{Dém. de 5.7.} L'implication b) $\Rightarrow$ c) est tautologique, compte tenu de la définition de la $\pi$-accessibilité de $M$, et de la prop. 5.3. \deleted{L'implication c) $\Rightarrow$ b) n'est autre que prop. \unsure{5.5.} (\unsure{revue|vue} et \unsure{corrigée}), que je viens de prouver. Donc $b\Leftrightarrow c$,} Et $a\Rightarrow b$ provient de la stabilité de $\mathrm{Filt}^{\pi}$ par $\underline{\lim}$ filtrantes de cardinal $\leq\pi$ (condition d)).

Il reste donc à prouver b) $\Rightarrow$ a) \inserted{(C'est ici qu'on doit utiliser l'hyp. particulière sur $\pi$)} \deleted{ce qui se réduit à prouver que b) implique que $M$ est stable par (petites) $\varinjlim$ filtrantes.} Une fois ceci vu, on saura l'équivalence de a,b,c, et \deleted{d'autre part cela implique aussitôt (via a et b)} d'autre part il est immédiat que a) $\Rightarrow$ c$'$), en

\scanpage{185}
\mspage{5\deleted{\unsure{8}}9}
% archive sheet number: 184

prenant $M_0=M\cap\mathrm{Filt}^{\unsure{\Pi|\pi}}$. D'autre part, $c'\Rightarrow c$, (\unsure{sans} hyp. sur $\pi$\unsure{?}), car \deleted{\unsure{donc}} \unsure{OPS|OBS} que la partie $M_0$ de $\mathrm{Filt}^{\Pi}$ est stable par $\underline{\lim}$ filtrantes de cardinal $\leq\pi$. Mais alors, connaissant l'équivalence de a,b,c, d'où \inserted{\unsure{aussitôt|aussi}} le corollaire, il \deleted{est} s'ensuit que l'enveloppe stable par $\underline{\lim}$ filtrantes de $M_0$ dans $\mathcal M$ \deleted{est} n'est autre que la partie accessible \deleted{\unsure{en}} déduite de $M_0$ dont il est question dans c, et on gagne.

Donc \deleted{\unsure{tout}} revient à prouver que, si tous les objets de $\mathrm{Filt}^{\Pi}$ sont accessibles, alors c implique a). Montrons d'abord que c) implique le critère d'appartenance \inserted{à $M$} (5.3.) (p. 51.), donc par la \deleted{démonstration \unsure{précédente} de} prop. 5.6. ``revue et corrigée'' en marge, il s'ensuit déjà que c) $\Rightarrow$ b), et il restera à établir b) $\Rightarrow$ a) i.e. que b) implique la stabilité de $M$ par $\underline{\lim}$ filtrantes quelconques (pas \unsure{seult} \unsure{gdes|grdes} devant $\pi$).

Que c) implique 5.3. va résulter du lemme suivant, appliqué à $I=\mathrm{Ss}_\pi(X)$ :

\uline{Lemme 5.8.} Soit $I$ un ensemble ordonné \uline{$\pi$-adapté}, i.e. filtrant et tel que toute partie filtrante $K$ de $I$, \inserted{de cardinal $\leq\pi$,} ait une borne supérieure $i_K$ dans $I$. \deleted{Considérons} \inserted{Je dis} \deleted{\unsure{l'ensemble des}} qu'une partie \deleted{filtrante} $J$ de $I$ est \deleted{strictement} \uline{$\pi$-cofinale} si elle est \deleted{filtrante} \inserted{cofinale} et si pour toute partie filtrante $K$ de $J$ de cardinal $\leq\pi$, la borne supérieure $i_K$ dans $I$ est

\scanpage{186}
\mspage{60}
% archive sheet number: 185

dans $J$. (\deleted{Il est immédiat que} \inserted{\unsure{Il s'ensuit|Il s'ensuit donc}} donc que $J$ est \unsure{lui-aussi|lui-même} $\unsure{\pi|\Pi}$-adapté, pour la str. d'ordre induite.) Ceci dit, \deleted{\illegible{1}} l'intersection de deux parties $\pi$-cofinales de $I$ est encore $\pi$-cofinale.

Il suffit de voir qu'elle est cofinale, car sa stabilité par bornes supérieures dans $I$ de parties filtrantes de cardinal est tautologique (si elle est vraie pour des $J_\alpha$, elle est vraie pour $\bigcap_\alpha J_\alpha$). \deleted{Quitte à remplacer} \inserted{Soit} $i_0\in I$, \deleted{$I$ par $I/i_0$, \unsure{on est ramené}} et construisons un élément de $J\cap J'$ qui \deleted{soit $\unsure{\geq|>} i_0$,} majore $i_0$. \deleted{Soit} On construit de proche en proche une \inserted{double} suite

\diagram{\begin{tikzcd}[ampersand replacement=\&,column sep=small,row sep=small]
 i_1\arrow[dr] \& \& i_2\arrow[dr,no head] \& \& \cdots\\
 \& i'_1\arrow[ur] \& \& i'_2\arrow[ur] \&
\end{tikzcd}}

avec les $i_\alpha$ dans $J$, les $i'_\alpha$ dans $J'$, et $i_1\geq i_0$. \deleted{Soit} On a donc

\[
\mathrm{Sup}\,i_\alpha=\mathrm{Sup}\,i'_\alpha\ \deleted{\unsure{\in I}},
\]

et si $i$ est la valeur commune des deux membres, on a $i\in J\cap J'$.

Ceci vu, pour prouver la validité de 5.3, on note que la partie $J'=\mathrm{Ss}^{(M^{\pi}_{\deleted{\unsure{0|\alpha}}})}_{\pi}(X)$ de $I=\mathrm{Ss}_\pi(X)$, formée des $Z_\alpha\subset X$ tels que $Z_\alpha\in M^\pi$, est stable par borne supérieure de parties filtrantes de cardinal $\leq\pi$, vu que $M^\pi$ est stable

\scanpage{187}
\mspage{61}
% archive sheet number: 186

par $\underline{\lim}$ de cardinal $\leq\unsure{\pi|\Pi}$. Donc si $J$ est une partie $\pi$-cofinale de $\mathrm{Ss}_{\pi}(X)$, $J\cap J'$ est cofinal dans $J$, i.e. dans $I$, sss $J'$ \deleted{\unsure{$\in$}} est \deleted{car alors elle sera} cofinal dans $I$ (i.e. sss $X\in M$, par la condition c). Car $J\cap J'$ cofinal $\Rightarrow J'$ cofinal tautologiquement, et $J'$ cofinal implique $J'$ $\pi$-cofinal, donc $J\cap J'$ $\pi$-cofinal par le lemme 5.8, donc cofinale.

Reste à prouver seulement que b) $\Rightarrow$ a), i.e. la stabilité de $M$ par limites inductives \inserted{filtrantes} quelconques, moyennant \deleted{le fait} b). Ici à nouveau, cela ne doit pas utiliser l'hypothèse particulière sur $\pi$, ni même l'hypothèse d'exactitude e) sur $\mathcal M$, mais être une histoire générale d'accessibilité :

\uline{Proposition} 5.9. Soit $\mathcal M$ une \deleted{partie} catégorie satisfaisant a) b) c) d) (p\unsure{ges} 44,45), $\pi\geq\pi_0$ un cardinal, $M$ une partie $\pi$-accessible de $\mathrm{Ob}\,\mathcal M$, i.e satisfaisant la condition \unsure{$1.b_{\pi})$} de \uline{$\overline{\mathrm{XIII}}$} p88. Alors $M$ est stable par $\varinjlim$ \unsure{filtrantes|filtrante} sss $M^\pi\overset{\mathrm{def}}=M\cap\mathrm{Filt}^{\pi}$ est stable par $\varinjlim$ filtrantes de cardinal $\leq\pi$.

Je n'ai pas envie d'expliciter ici la démonstration, ayant l'impression pénible de refaire à \unsure{perpète} la même démonstration (ici et dans \uline{$\overline{\mathrm{XIII}}$}). Je \unsure{devrai|devrais} revenir là-dessus dans une \unsure{rédaction|situation} plus décantée.

\scanpage{188}
\mspage{62}
% archive sheet number: 187

Je vais quand même esquisser la démonstration. Soit $X=\varinjlim_I X_i$, $I$ filtrant, les $X_i$ dans $M$, il faut prouver que $X\in M$ (moyennant la \inserted{l'hyp. de} stabilité sur $M^{\unsure{\pi|\Pi}}=M\cap\mathrm{Filt}^{\Pi}$).

\diagram{\begin{tikzcd}[ampersand replacement=\&,column sep=small,row sep=small]
 \& C\arrow[r]\arrow[d]\arrow[dl,dashed,"{\unsure{z|3|\mathfrak z}}"'] \& I\arrow[d,"{\unsure{\varphi|\phi}}"]\\
 \underline M^\pi \& \underline M^\pi/\underline M\arrow[l,dashed,"{\text{``source''}}"]\arrow[r,"{\text{``but''}}"'] \& \underline M
 \arrow[ul,phantom,"{\text{cart}}" description]
\end{tikzcd}}

$C$ est une catégorie cofibrée sur $I$, dont les fibres $C_i$ sont

\[
 C_i\simeq\underline M^\pi/X_i,
\]

or comme $X_{\unsure{i|0}}\in M$, on voit, à cause de la condition de $\pi$-accessibilité, que $M^\pi/X_i$ est \uline{filtrante}. \deleted{\unsure{Cela} En effet} \inserted{\unsure{Car ce fait|Car ceci}} \deleted{\unsure{Lemme} \unsure{de ce fait}} \inserted{dégageons un}

\margin{\rule{0.4pt}{1.4cm}}\uline{Lemme} 5.\deleted{\unsure{9}}10. Soit $M$ une partie $\pi$-accessible de $\mathcal M$. Alors pour tt $\deleted{\unsure{Z}}Y\in\deleted{\unsure{X}}M$, $M^\pi/\deleted{\unsure{Z}}Y$ est une catégorie \uline{filtrante}, et $\mathrm{Ss}_\pi(Y)$ y est cofinale.

\deleted{En effet, les $Z\longrightarrow Y$ avec $Z$ dans $M^\pi$ (ou seulement dans $\mathrm{Filt}^{\pi}$) se factorisent par $Z\longrightarrow Z_\alpha\hookrightarrow Y$, avec $Z_\alpha$ dans $\mathrm{Ss}_\pi(Y)$, or \deleted{\unsure{les}} si $Y\in M$, les $Z_\alpha$ tels que $Z_\alpha\in M$ sont cofinaux, donc les $Z\longrightarrow Y$ ($Z\in M^\pi$) se factorisent par $Z_\alpha\lhook\joinrel\longrightarrow Y$ avec $Z_\alpha\in M^\pi$. Cela montre que si $Z\to Y,Z'\to Y$ sont dans $M^\pi/Y$, deux obj sont majorés par un $Z_\alpha\to Y$ $\in\mathrm{Ss}_\pi(Y)\unsure{/Y}\subset M^\pi/Y$. D'autre}

\scanpage{189}
\mspage{63}
% archive sheet number: 188

\deleted{part, si $\deleted{\unsure{X}}Z\longrightarrow Y,Z'\unsure{-|\longrightarrow}Y$ sont dans $M^{\unsure{\pi|\Pi}}/Y$, et si on a une double flèche}

\deleted{\diagram{\begin{tikzcd}[ampersand replacement=\&,column sep=small,row sep=small]
 \deleted{Z}\arrow[rr,shift left=.6ex]\arrow[rr,shift right=.6ex]\arrow[dr] \& \& \deleted{Z'}\arrow[dl]\\
 \& \deleted{Y} \&
\end{tikzcd}}}

\deleted{en factorisant par $Z'\longrightarrow\underset{\substack{\rotatebox{-90}{$\in$}\\\deleted{\unsure{M}}\mathrm{Ss}^{\pi}_{\pi}(Y)\subset M^\pi/Y}}{Z_\alpha}\lhook\joinrel\longrightarrow Y$, on trouve $Z'\longrightarrow Z_\alpha$ qui égalise la double flèche.}

Considérons le foncteur canonique

\[
 I=\mathrm{Ss}_{\pi}(Y)\xrightarrow{\unsure{\varphi|\phi}}\underline M^\pi/Y=\underset{\unsure{9|g|\varphi}}{I'},
\]

où la catégorie source \inserted{$I$} est filtrante. Par SGA 4 I 8.1.3, ce foncteur est cofinal sss il satisfait deux conditions \unsure{F1, F2|F1, F7}, et alors \deleted{\unsure{$I$}} $I'$ est également filtrante. La \deleted{\unsure{même}} condition F1, \deleted{\unsure{dit que}} est la condition habituelle que tt $i'\in\mathrm{Ob}\,I'$ est ``majoré'' par un $\unsure{\varphi(i)}$, i.e. \unsure{$\forall$} $Z\longrightarrow Y$, avec $Z\in\underline M^\pi$, est ``majoré'' par un \inserted{\uline{mono}} $Z_\alpha\hookrightarrow Y$, \inserted{avec $Z_\alpha\in M^\pi$} ce qui résulte du fait que \deleted{si $Z$} c'est majoré par $Z_\alpha\lhook\joinrel\longrightarrow Y$ avec $Z_\alpha\in\mathrm{Filt}^{\pi}$, \deleted{et que} i.e $Z_\alpha\in\mathrm{Ss}_\pi(Y)$, et que parmi ces $Z_\alpha$ ceux qui sont dans $M^\pi$ sont cofinaux lorsque $Y\in M$.

% slanted
\margin{Comme $I$ est filtrante}

L'autre condition \unsure{*} est qu'une double flèche $i'\mathrel{\substack{\xrightarrow{f',g'}\\[-.5ex]\longrightarrow}}\varphi(i)$ \unsure{puisse|pour}

\scanpage{190}
\mspage{64}
% archive sheet number: 189

s'égaliser par une flèche $\unsure{\varphi|\phi}(h)$, avec $h:i\longrightarrow j$ \deleted{\unsure{pour|dans}} Mais en fait \deleted{\illegible{2}} pour un $i'\in\mathrm{Ob}\,I'$ i.e un $Z\longrightarrow Y$ donné ($Z\in M^{\unsure{\pi|\Pi}})$ et un $i\in\mathrm{Ob}\,I$ i.e. un $Z_\alpha\lhook\joinrel\longrightarrow Y$ donné ($Z_\alpha\in\mathcal M^\pi$), il existe au plus \uline{une} flèche de \deleted{\unsure{$I'$}, $I=\unsure{M}$} dans $I'=\underline M^\pi/Y$ de $Z\longrightarrow Y$ dans $Z_\alpha\lhook\joinrel\longrightarrow Y$. Donc la condition F2 est trivialement satisfaite (on prend \deleted{\unsure{$j$}} $h=\mathrm{id}_i$).

\uline{Corollaire 5.11.} Si $(X_i)_{\unsure{i'|i}\in I}$ est un syst. inductif filtrant, avec les $X_i$ dans $M^\pi$, alors la catégorie $C=\underline M^\pi/M\times_{\underline{M}}(I,\varphi)$ construite p. 62 est filtrante, et le foncteur canonique \unsure{(fibrant)} $C\longrightarrow I$ est cofinal.

\deleted{Cela se} Que $C$ soit filtrante \inserted{et $C\to I$ cofinal,} résulte \unsure{immédiatement|également} du fait que $I$ l'est, \deleted{et} que les fibres de $C$ le sont, et que $C$ est cofibrée sur $I$. (Mérite un lemme, à \deleted{\unsure{toutes}} fins utiles.) \deleted{\unsure{Donc}}

\deleted{\unsure{On a maintenant}} \uline{NB} On peut \unsure{remplacer|compléter} $C$ par $C_0$ dont les fibres sont les $\mathrm{Ss}^{M^\pi}_{\pi}(X_i)\subset\underline M^{\pi}$; \unsure{à partir de|cofinale de} $C_i=\underline M^\pi/X_i$,

\scanpage{191}
\mspage{65}
% archive sheet number: 190

on trouve une ss-catégorie de $C$ \unsure{cofinal|cofinale} dans $C$, et qui est cette fois ordonnée \deleted{\unsure{tout}} comme $I$.

On a

\[
X_i=\underline{\lim}_{Z_\alpha\in C_i}\deleted{\unsure{C_i|C}}Z_\alpha,
\]

\deleted{donc} donc

\[
X=\underline{\lim}_{\alpha\in C}Z_\alpha,
\]

au total on est arrivé à représenter $X$ comme $\underline{\lim}$ \deleted{\unsure{filtrante}} d'éléments $Z_\alpha$ qui sont (non seulement dans $M$, mais $=$ \deleted{\unsure{dans}}) dans $M^{\unsure{\pi|\Pi}}$.

Changeons de \unsure{notation|notations}, en prenant

\[
\begin{aligned}
 X&=\varinjlim_{i\in I}X_i,\qquad I\ \text{un ordonné filtrant}\\
 &\hspace{5em}\qquad X_i\in M^\pi\quad\forall i\in I,
\end{aligned}
\]

soit $I'$ l'ens \deleted{\unsure{ordonné}} des parties \uline{filtrantes} $J\subset I$ telles que $\mathrm{card}\,J\leq\pi$, $I'$ ordonné par inclusion. Alors $I'$ est grand devant $\pi$, et on a un système inductif

\[
\deleted{\unsure{\varinjlim}}\ (X_J)_{J\in I'},\qquad X_J=\varinjlim_{i\in J}X_i.
\]

Par l'hypothèse de stabilité de $M^\pi$ (qui n'a pas encore servi), les $X_J$ sont dans $M^\pi$. On a donc

\[
X=\varinjlim_{J\in I'}X_J \qquad
\begin{aligned}
 &\text{les $X_J$ dans $M^\pi$}\\
 &\text{$I'$ grd devant $\pi$}
\end{aligned}
\]

donc $X\in M$ puisque $M$ $\pi$-admissible, qed.

\scanpage{192}
\mspage{66}
% archive sheet number: 191

\uline{Définition} 5.12. 1) Soit $\mathcal M$ une catégorie satisfaisant les conditions a) b) c) d) (p. 44.\unsure{--|$=$}45), \inserted{soit $\unsure{\pi|\Pi}\geq\pi_0$ un card\unsure{inal}.} Une partie $M$ de $\mathrm{Ob}\,\mathcal M$ est dite \uline{$L$-$\pi$-accessible}, si elle est $\pi$-accessible et stable par \inserted{\uline{petites}} $\underline{\lim}$ filtrantes\unsure{*}, ou ce qui revient au $=$ \inserted{(5.9)}, si elle est $\pi$-accessible et si $M^\pi=M\cap\mathrm{Filt}^{\Pi}$ est stable par $\underline{\lim}$ filtrantes de cardinal $\leq\pi$.

% slanted
\margin{On dit alors que $M$ est \uline{$L$-stable}}

(\uline{NB} Si les objets de $\mathrm{Filt}^{\Pi}$ sont $\pi$-accessibles, 5.8 nous dit qu'il cela signifie aussi qu'il \inserted{suffit pour cela qu'il} existe une partie $M_0$ de $\mathrm{Filt}^{\Pi}$ telle que $M$ soit la \inserted{$L$-}enveloppe (enveloppe stable par $\underline{\lim}$ filtrantes) de $M_0$ -- condition qui est de \unsure{ttes|tte} \unsure{façons|façon} tjrs nécessaire, en prenant $M_0=M^\pi$.)

2) Soit $M$ une partie de $\mathrm{Ob}\,\mathcal M$. On \deleted{\unsure{dit elle}} qu'elle est \uline{$L$-accessible}, si elle satisfait à l'une des \deleted{deux} \inserted{trois} conditions équivalentes suivantes :

a) Il existe un cardinal $\pi$ tel que $M$ soit $L$-$\pi$-accessible, \deleted{\unsure{b).}}

\inserted{a$'$) $M$ est accessible et $L$-stable.}

b) Il existe une petite partie $M_0$ de $\mathrm{Ob}\,\mathcal M$, telle que $M$ soit la $L$-enveloppe de $M_0$.

% slanted
\margin{$\longleftarrow$ Cette condition ne dépend pas de la filtration cardinale choisie de $\mathcal M$}

b$'$) Il existe une partie $M_0$ de $\mathcal M$ contenue dans une partie $\mathrm{Filt}^{\pi}$ (ou ce qui revient au $=$, telle que la \inserted{sous} catégorie pleine \deleted{\unsure{de}} $\underline M_0$ de $\mathcal M$ soit \unsure{ess\textsuperscript{t.}} petite, i.e. égale à une petite catégorie), telle que $M$ soit la $L$-enveloppe de $M_0$.

\scanpage{193}
\mspage{67}
% archive sheet number: 192

L'équivalence de a) et a$'$), \unsure{ou|resp.} de b) et b$'$), est $\pm$ tautologique. L'équivalence de a) et b) résulte aussitôt du NB inclus au \uline{1°}, compte tenu que les $\unsure{\pi|\Pi}$ tels que les objets de $\mathrm{Filt}^{\Pi}$ soient $\pi$-accessibles sont cofinaux dans l'ens des cardinaux $\in\unsure{\mathcal U|\mathscr U|U}$. (Cf 3.2, p. 16)

\margin{\rule{0.4pt}{3cm}}\uline{Proposition 5.13.} Toute réunion finie de parties $L$-$\pi$-accessibles de $\mathrm{Ob}\,\mathcal M$ est $L$-$\pi$-accessible. L'intersection d'une famille $(M_\alpha)_{\alpha\in A}$ de parties $L$-$\pi$-accessibles, avec $\mathrm{card}\,A\leq\pi$, est $L$-$\pi$-accessible.

\uline{Dém} La première assertion provient du fait qu'une réunion \uline{finie} de parties $L$-stables de $M$ est $L$-stable, \deleted{ce qui} et qu'on a \unsure{itou|idem} pour la notion de $\pi$-accessibilité. Les deux choses se prouvent d'ailleurs pratiquement de la même façon, en utilisant le

\uline{Lemme} 5.14. Soit $I$ un ens. ordonné filtrant, $(I_\alpha)_{\alpha\in A}$ une famille \uline{finie} \inserted{non vide} de parties de $I$, telle que $\bigcup I_\alpha=I$. Alors il existe $\alpha\in A$ tel que $I_\alpha$ soit cofinal, \deleted{dans $I$ et $=\ \alpha\in A$ et $i_0\in I$ tels que $I_\alpha\supset I/i_0$}

\scanpage{194}
\mspage{68}
% archive sheet number: 193

\deleted{Par réc\unsure{.} sur $\mathrm{card}\,A$, on est ramené au cas de \uline{deux} parties $I_1,I_2$. Supposons que $I_1$ n'est pas cofinal, prouvons que $I_2$ est cofinal. Soit $i_0\in I$, \deleted{\unsure{tel que}} tel qu'il n'existe pas d'élément de $I_1$ $\geq i_0$. On a donc $I_{/i_0}\subset I\setminus I_1$, donc $I_{/i_0}\subset I_2$ puisque $I_1\cup I_2=I$, donc $I_2$ est cofinal, $I_{/i_0}$ l'étant. \uline{NB} Cet argument montre $=$ mieux, savoir qu'il existe un $\alpha$ \inserted{et un $i_0\in I$} tel que $I_\alpha\supset I_{/i_0}$.}

Dém par l'absurde. Si pour \deleted{tt $\alpha$,} aucun $\alpha$, $I_\alpha$ \unsure{n'était|n'étant} cofinal, on en conclurait que pour tt $\alpha$ il existe un $i_\alpha$, tel que $I_{/i_\alpha}\cap I_\alpha=\varnothing$. Soit $i$ \deleted{tel que} majorant les $i_\alpha$, on aura donc $I_{/\deleted{\unsure{\alpha|i_\alpha}}i}\cap I_\alpha=\varnothing\quad\forall\alpha\in A$, contrairement à l'hypothèse que $\bigcup I_\alpha=I$.

Le cas d'une \uline{intersection} \deleted{\unsure{finie de p} \unsure{résulte|arbitraire} du lemme la définition b) \unsure{sous la} forme} de parties $L$-$\unsure{\pi|\Pi}$-accessibles : elle est bien sûr $L$-stable, donc il reste à \unsure{montrer|prouver} qu'elle est $\pi$-accessible, i.e. satisfait la condition $b_\pi$ de loc. cit.

\scanpage{195}
\mspage{69}
% archive sheet number: 194

Soit donc \deleted{$X=$}

\[
 X=\varinjlim_{i\in I} X_i \qquad X_i\in\mathrm{Filt}^{\unsure{\Pi|\pi}}
\]

une représentation $\pi$-adaptée de $X$, avec $I$ un ens. ordonné $\pi$-adapté. Donc on a

\[
\begin{aligned}
 X\in M_\alpha \quad\Longleftrightarrow\quad
 &\text{l'ens $I_{M_\alpha}$ des $\deleted{\unsure{\alpha}}i\in I$ tels que}\\
 &X_{\deleted{\unsure{\alpha}}i}\in M_\alpha\ \text{est cofinal}\\
 &\text{\unsure{en|dans} $I$,}
\end{aligned}\tag*{\unsure{(*)}}
\]

D'ailleurs, comme $\underline M^{\deleted{\unsure{\alpha}}\pi}_{\alpha}=M_{\unsure{0|\alpha}}\cap\mathrm{Filt}^\pi$ est stable par $\underline{\lim}$ filtrantes de cardinal $\leq\pi$, \unsure{l'on|on} voit \deleted{\unsure{tout}} de suite que la partie $I_{M_\alpha}$ de $I$ est stable par bornes supérieures dans $I$ des parties $J$ de $I_{M_\alpha}$ filtrantes et de cardinal $\leq\pi$. Donc $I_{M_\alpha}$ cofinal équivaut à \uline{$I_{M_\alpha}$ $\pi$-cofinal}. (cf 5.8 page 59). Soit $M=\bigcap_{\alpha\in A}M_\alpha$, je veux prouver que

\[
\begin{aligned}
 X\in M \quad\Longleftrightarrow\quad &\text{l'ens $I_M$ ($=\bigcap I_{M_\alpha}$)}\\
 &\text{est \deleted{\unsure{adapté}} cofinal dans}\\
 &I
\end{aligned}\tag*{\unsure{(**)}}
\]

Or \unsure{(*)} implique que $X\in M$ ($=\bigcap M_\alpha$) revient à : $I_{M_\alpha}$ $\pi$-cofinal pour \unsure{tt} $\alpha$. Donc \unsure{au total|en tout} on veut prouver

\[
\begin{aligned}
 &I_{M_\alpha}\,\deleted{\unsure{\pi}}\text{-cofinal}
 &&\Longleftrightarrow\quad \bigcap I_{M_\alpha}\ \text{est $\deleted{\unsure{\pi}}$-cofinal}\\
 &\text{dans $I$ pour tt} &&\hspace{5em}\text{dans }I.\\
 &\alpha
\end{aligned}\tag*{\unsure{(***)}}
\]

